% Vietnamese translation for OI Public Library
% Source-PDF: 其他 OTHERS/NOI2016-旷野大计算-题解.pdf
\documentclass[11pt]{article}
\usepackage{oipl-vn}

\usepackage{array}

\title{Lời giải bài ``Đại tính toán nơi hoang dã''}
\author{Lu Kaifeng\\Người kiểm đề: Wang Yisong, Yu Jiping, Zhang Ruizhe}
\date{26 tháng 7 năm 2016}

\begin{document}
\maketitle

\origtitle{NOI2016 Wilderness Computation solution}
\sourcepath{Others / NOI2016 Wilderness Computation solution.pdf}

\begin{abstract}
Tài liệu gốc là bộ slide lời giải cho một bài dựng mạng tính toán. Mỗi test
yêu cầu ghép các nút tính cơ bản thành một mạng càng ít nút càng tốt. Bản
dịch này chuyển nội dung slide thành bài viết, giữ lại các công thức, bảng
điểm, ý tưởng chấm điểm và số thao tác; các slide chỉ có hình minh họa được
kiểm tra trực quan và gộp vào prose khi chúng không mang thêm dữ liệu thuật
toán.
\end{abstract}

\tableofcontents

\section{Tán gẫu}

Xin chào mọi người! Tôi, vfk, đã quay lại.

Lần này tôi ra cho mọi người một bài ``đại tính toán nơi hoang dã'', không
biết mọi người có kiếm được thật nhiều điểm nơi hoang dã không?

Nghe nói Vương quốc Bọ Chét đã chiếm đề bài NOI rồi, nhưng thật ra chúng tôi
viết hai đề riêng biệt. 233.

\section{Khái quát bài toán}

Hai slide đầu của phần này dùng hình minh họa: một cảnh trong phim
\textit{Red Cliff}, rồi một sơ đồ mạng nơ-ron truyền thẳng nhỏ. Nội dung cần
giữ lại là bài toán mô hình hóa việc tính toán bằng một mạng có hướng gồm
các nút thao tác.

Các loại nút trong mạng tính toán như sau.

\begin{center}
\small
\begin{tabular}{|>{\raggedright\arraybackslash}p{0.20\textwidth}|c|>{\centering\arraybackslash}p{0.28\textwidth}|>{\raggedright\arraybackslash}p{0.28\textwidth}|}
\hline
Tên nút & Ký hiệu & Kết quả tính & Đặc điểm \\
\hline
Nút nhập & \texttt{I} & đầu vào & Nhập và xuất cơ bản \\
\hline
Nút xuất & \texttt{O} & đầu ra & Nhập và xuất cơ bản \\
\hline
Nút cộng & \texttt{+} & $x_t=x_i+x_j$ & Cộng trừ \\
\hline
Nút dịch hằng & \texttt{C} & $x_t=x_i+c$ & Cộng trừ \\
\hline
Nút đổi dấu & \texttt{-} & $x_t=-x_i$ & Cộng trừ \\
\hline
Nút dịch trái & \texttt{<} & $x_t=x_i\cdot 2^k$ & Nhân chia với lũy thừa của $2$ \\
\hline
Nút dịch phải & \texttt{>} & $x_t=x_i/2^k$ & Nhân chia với lũy thừa của $2$ \\
\hline
Nút kiểu S & \texttt{S} & $x_t=s(x_i)$ & Một thao tác phi tuyến \\
\hline
Nút so sánh & \texttt{P} & $x_t=\operatorname{cmp}(x_i,x_j)$ & Mạnh nhưng bị trừ điểm \\
\hline
Nút Max & \texttt{M} & $x_t=\max(x_i,x_j)$ & Mạnh nhưng bị trừ điểm \\
\hline
Nút nhân & \texttt{*} & $x_t=x_i\cdot x_j$ & Mạnh nhưng bị trừ điểm \\
\hline
\end{tabular}
\end{center}

Hàm sigmoid được định nghĩa bởi
\[
  s(x)=\frac{1}{1+e^{-x}}.
\]

Đề cho $10$ nhiệm vụ tính toán. Với mỗi nhiệm vụ, ta dùng một số thao tác cơ
bản để dựng một mạng tính toán. Số nút dùng càng ít càng tốt.

\section{Điểm số}

Điểm cao nhất và điểm trung bình theo từng test như sau.

\begin{center}
\begin{tabular}{|c|c|c|}
\hline
Test & Điểm cao nhất & Điểm trung bình \\
\hline
1 & 10 & 9.7 \\
2 & 10 & 9.6 \\
3 & 10 & 8.4 \\
4 & 9 & 5.6 \\
5 & 10 & 9.1 \\
6 & 9 & 2.3 \\
7 & 10 & 1.3 \\
8 & 10 & 4.5 \\
9 & 10 & 1.2 \\
10 & 10 & 0.4 \\
\hline
\end{tabular}
\end{center}

Thống kê chung:
\begin{itemize}
  \item điểm cao nhất: $97$ điểm, của Yuan Yutao;
  \item điểm trung bình: $52.1$ điểm;
  \item trung vị: $52$ điểm;
  \item mốt: $52$ điểm.
\end{itemize}

Hai slide ``bí mật lớn của điểm số nơi hoang dã'' và slide ``lời phàn nàn nơi
hoang dã'' trong bản gốc không có nội dung văn bản trích xuất ngoài tiêu đề.

\section{Test 1: làm quen thao tác cơ bản}

Nhập $a,b$, tính
\[
  -2a-2b.
\]

\subsection*{Cách $5$ điểm}

Dùng biểu thức
\[
  -(a\ll 1)+\bigl(-(b\ll 1)\bigr).
\]
Tổng cộng $8$ thao tác.

\subsection*{Cách $10$ điểm}

Dùng biểu thức
\[
  -\bigl((a+b)\ll 1\bigr).
\]
Tổng cộng $6$ thao tác.

\section{Test 2: làm quen thao tác cơ bản}

Nhập $a$, tính
\[
  \frac{1}{1+e^{17a}}.
\]

\subsection*{Cách $4$ điểm}

Dùng
\[
  s\bigl(-(a+a+a+\cdots+a)\bigr).
\]
Tổng cộng $20$ thao tác.

\subsection*{Cách $10$ điểm}

Dùng
\[
  s\bigl(-(a+(a\ll 4))\bigr).
\]
Tổng cộng $6$ thao tác.

\section{Test 3: xét dấu}

Nhập $a$, tính
\[
  \operatorname{sgn}(a)=
  \begin{cases}
    -1, & a<0,\\
    0, & a=0,\\
    1, & a>0.
  \end{cases}
\]

\subsection*{Cách $6$ điểm}

Dùng thao tác bị trừ điểm. Trước hết dùng $a\gg 1000$ để dựng ra $0$, rồi đưa
$0$ và $a$ vào nút so sánh.

\subsection*{Cách $10$ điểm}

Đặt
\[
  s(x\ll 500)\approx p(x)=
  \begin{cases}
    0, & x<0,\\
    1/2, & x=0,\\
    1, & x>0.
  \end{cases}
\]
Chỉ cần tịnh tiến và co giãn tuyến tính $p(a)$ là nhận được
$\operatorname{sgn}(a)$.

\section{Test 4: trị tuyệt đối}

Nhập $a$, tính $\lvert a\rvert$.

\subsection*{Cách $2$ điểm}

Dùng thao tác bị trừ điểm: dùng nút so sánh để tính
$\operatorname{sgn}(a)$, sau đó dùng nút nhân để tính
\[
  \lvert a\rvert=a\cdot \operatorname{sgn}(a).
\]

\subsection*{Cách $6$ điểm}

Dùng ý tưởng của test 3 để tính $\operatorname{sgn}(a)$, rồi dùng nút nhân
\[
  \lvert a\rvert=a\cdot \operatorname{sgn}(a).
\]
Hoặc trực tiếp dùng nút Max cũng tính được trị tuyệt đối.

\subsection*{Cách $10$ điểm}

Thử nghĩ: nếu bọc thêm một lớp $s$ thì sao? Với một giá trị $x$, xét
\[
  s\bigl(x+(p(x)\ll 500)\bigr)\approx
  \begin{cases}
    s(x), & x<0,\\
    1, & x\ge 0.
  \end{cases}
\]
Nếu vế đầu không phải $s(x)$ mà là chính $x$, mọi chuyện sẽ rất hay.

Khi đại lượng là vô cùng bé thì điều gì xảy ra? Theo định nghĩa đạo hàm, khi
$x$ rất gần $0$,
\[
  \frac{s(x)-s(0)}{x-0}\approx s'(0)=\frac14.
\]
Vì vậy, nếu $x$ rất nhỏ, ta có thể suy ngược $x$ từ $s(x)$.

Cụ thể,
\[
  s\bigl((x\gg 150)+(p(x)\ll 500)\bigr)\approx
  \begin{cases}
    1/2+(x\gg 154), & x<0,\\
    1, & x\ge 0.
  \end{cases}
\]
Từ đó dựng được $\min(x,0)$, và do đó dựng được $\lvert x\rvert$.

Sau một số chỉnh hằng số cần thiết, có thể lấy $10$ điểm. Mã giả:

\begin{verbatim}
x = in();
p = s((x + 1e-30) << 500) << 152;
y = s((x >> 150) + p);
r = x + ((-y + 0.5) << 153) + p;
out(r);
\end{verbatim}

Tổng cộng $14$ thao tác.

\section{Test 5: chuyển đổi nhị phân}

Xuất số nguyên nhị phân được biểu diễn bởi
\[
  a_1,\ldots,a_{32}.
\]

\subsection*{Cách $10$ điểm}

Trực tiếp dùng cộng và nhân là đủ, tổng cộng $95$ thao tác.

\section{Test 6: chuyển đổi nhị phân}

Xuất biểu diễn nhị phân của số nguyên
\[
  a\in [0,2^{32}).
\]

Ý tưởng cơ bản:
\begin{verbatim}
x = in();
for (i = 31; i >= 0; i--)
{
    if (x >= (1 << i))
        b[i] = 1;
    else
        b[i] = 0;
    x = x + (-(b[i] << i));
}
out(b);
\end{verbatim}

\subsection*{Cách $6$ điểm}

Dùng nút so sánh là có thể qua.

\subsection*{Cách $9$ điểm}

Dùng hàm $p(x)$ đã nói ở trên để hiện thực nút so sánh, kèm một số chỉnh hằng
số cần thiết. Chẳng hạn có thể dùng $x\ll 500$ thay cho $x$ trong tính toán,
vì khi tính hàm $p$ đằng nào cũng phải dịch trái.

\subsection*{Cách $10$ điểm}

Chú ý rằng khi $i=0$, ta có thể đặt thẳng $b[0]=x$. Tổng cộng $190$ thao tác.

\section{Sáu test đầu}

Sáu test đầu đều là những thao tác rất cơ bản. Khi đã chơi thật rõ các thao
tác cơ bản này, bốn test thuật toán nâng cao phía sau cũng không còn khó nữa.
Hơn nữa các tham số chấm điểm phía sau được đặt rất thoáng.

\section{Test 7: phép toán nhị phân}

Nhập $a,b$, tính $a\operatorname{ xor }b$.

Trước hết cần làm tốt test 5 và test 6. Sau khi chuyển đổi nhị phân, ta chỉ
cần xét phép xor của hai biến $01$ là $x$ và $y$.

\subsection*{Cách $10$ điểm}

\begin{verbatim}
s = x + y;
return s + (-p(s + (-1.5)) << 1);
\end{verbatim}

Tổng cộng $603$ thao tác. Người kiểm đề saffah tối ưu được xuống $539$ thao
tác.

\section{Test 8: chia cho hằng số tổng quát}

Nhập $a$, tính $a/10$.

\subsection*{Cách $6$ điểm}

Dùng thao tác bị trừ điểm. Trước hết dùng $x\gg 1000$ để dựng ra $0$, rồi
cộng thêm $0.1$ để nhận $0.1$, sau đó dùng nút nhân để tính $0.1x$.

\subsection*{Cách chưa rõ điểm}

Có thể nhân $a$ với $2^{32}$, chuyển thành số nguyên, rồi viết một phép chia
nhị phân.

\subsection*{Cách $7$ đến $8$ điểm}

Viết $0.1$ dưới dạng nhị phân, biến nó thành tổng của một số hạng $x$ dịch
phải.

\subsection*{Cách $9$ điểm}

Trên nền cách trước, dùng phép trừ để chỉnh hằng số, có thể bỏ đi một đoạn
liên tiếp các bit $1$ trong biểu diễn nhị phân. Hoặc dùng tuyệt chiêu ``vẻ đẹp
của chu kỳ'': tính ra một chu kỳ rồi nhân đôi.

\subsection*{Cách $10$ điểm}

Theo định nghĩa đạo hàm, khi $x$ rất gần $x_0$,
\[
  \frac{s(x)-s(x_0)}{x-x_0}\approx s'(x_0).
\]
Vì vậy chỉ cần giải một $x_0$ sao cho
\[
  s'(x_0)=0.1.
\]
Cách này chỉ cần $7$ thao tác. Do bài chỉ cần độ chính xác $10^{-9}$, có thể
dùng \texttt{double} để tính $x_0$, rồi đưa vào checker để tính
$s(x_0)$ với độ chính xác cao.

\section{Test 9: sắp xếp}

Nhập $a_1,\ldots,a_{16}$ và sắp xếp chúng.

Nhiệm vụ tính toán này cho thấy rõ giới hạn của mô hình mạng tính toán: chỉ
những thuật toán như sắp xếp nổi bọt hoặc sắp xếp chọn mới thích hợp để chạy.
Dĩ nhiên có các thuật toán độ phức tạp thấp hơn như bitonic sort, nhưng trong
bài này dùng sắp xếp nổi bọt đã đủ.

Điểm mấu chốt là hiện thực bộ so sánh. Ví dụ:

\subsection*{Cách $10$ điểm}

\begin{verbatim}
s = x[i] + x[j];
x[i] = x[i] + min(-x[i] + x[j], 0);
x[j] = s + -x[i];
\end{verbatim}

Tổng cộng $2072$ thao tác. Cách hiện thực cụ thể của $\min(x,0)$ có thể tham
khảo test 4.

\section{Test 10: phép nhân theo modulo}

Nhập $a,b,m$, tính
\[
  a\cdot b\bmod m.
\]

\subsection*{Cách $10$ điểm}

Hiện thực một phép ``cộng nhanh'', tức lấy thuật toán lũy thừa nhanh rồi thay
phép nhân trong đó bằng phép cộng. Cần chỉnh hằng số một chút.

Ta chỉ cần hiện thực phép lấy modulo $m$ cho một số trong đoạn $[0,2m)$.

Cách này đã đủ điểm tối đa. Nhưng có thể dùng ít thao tác hơn nữa không?

\subsection*{Cách $10$ điểm mạnh hơn}

Theo khai triển Taylor,
\[
  s(x+\ln 3)
  = \frac14+\frac{3}{16}x+\frac{3}{64}x^2
    -\frac{1}{256}x^3-\frac{5}{1024}x^4+O(x^5).
\]
Từ đó có thể dựng $x^2$, rồi dựng tiếp $xy$. Sau đó mô phỏng ý tưởng chuyển
số nguyên sang nhị phân để hiện thực phép lấy modulo là xong.

Có thể làm chỉ với $819$ thao tác.

\section{Bài học}

\subsection*{Mô-đun hóa}

Slide gốc dùng một ảnh minh họa cho ý tưởng mô-đun hóa. Trong lời giải, tinh
thần này thể hiện ở việc tách các khối nhỏ như xét dấu, dựng $\min(x,0)$,
chuyển nhị phân, bộ so sánh và lấy modulo, rồi tái sử dụng chúng cho các test
sau.

\subsection*{Rời rạc và liên tục đều có thế mạnh}

Các test sau vừa dùng thủ thuật rời rạc như biểu diễn nhị phân, xor, sắp xếp
nổi bọt và cộng nhanh, vừa dùng xấp xỉ liên tục qua sigmoid, đạo hàm và khai
triển Taylor. Hai hướng này bổ sung cho nhau.

\subsection*{Mạng nơ-ron}

Slide gốc kết thúc phần bài học bằng các ảnh minh họa về mạng nơ-ron. Điều
này tương ứng với mô hình của đề: một mạng gồm các nút tính toán, trong đó
nút sigmoid đóng vai trò thao tác phi tuyến quan trọng.

\section{Kết thúc}

Hoàn thành một thành tựu trong đời: ra đề ở NOI.

Cảm ơn CCF đã cung cấp nền tảng học tập và giao lưu. Cảm ơn Wang Yisong và Yu
Jiping đã kiểm đề và siết tham số chấm. Cảm ơn Zhang Dizhu đã cân nhắc câu
chữ của đề bài.

Cảm ơn mọi người!

\end{document}
